A list of the Linux commands referred to in this paper is included in \autoref{AppendixLinuxCommands}. This list is not an exhaustive list of all commands available on Linux-based distributions, it is based on the commands found on \textcite{CommandLineReference}, which includes most if not all the built-in bash commands. For a full list of Linux commands, see \textcite{LinuxManPages}. 

To begin with, we can derive that commands on the Linux commandline can affect most of anything on the system. Processes, files, network connections, users and even hardware can be affected by commands. This means that in order to understand what can be done locally, we need to understand how the system works at its most fundamental level and what commands give way to manipulate it.

\subsection{Processes and Jobs}
\label{ProcessesAndJobs}

Processes are a fundamental unit of execution in the Linux operating system, the management of which oversees all actions taken on the system and by the system. Understanding how these processes work, how they can be monitored and affected by the user, commands, and other processes is crucial in understanding how AI agents could operate in our simulated environment, and how they could potentially interact with each other.

For references on processes we refer to our previous work \citetitle{AsyncSync} \parencite{AsyncSync}, this report will not cover the finer details of process or thread management, but rather what they mean for the commandline.

\subsubsection{Process Architecture and Concepts}
\label{ProcessArchitecture}

There are four important keywords in this category, which are \textit{processes}, \textit{threads}, \textit{pipelines} and \textit{jobs}. Processes can be somewhat loosely defined as we discovered in our prior study \parencite{AsyncSync}, however they are generally considered to be instances of code that are running in an isolated memory space. Each process has its own virtual address space, preventing it from directly accessing the memory of other processes and this isolation is enforced by the kernel and provides fundamental security boundaries between processes.

When a process is created through the \texttt{fork()} system call, it receives a unique Process ID (PID). The newly created child process is initially an exact copy of the parent process, including its memory space, open file descriptors, and execution context. The child can then use \texttt{exec()} to replace its memory space with a new program. This fork-exec model is fundamental to how shells launch commands and how processes spawn subprocesses.

Threads can be generally split into two categories, user-level threads and kernel threads, where the former are managed by the process itself and the latter are managed by the operating system's scheduler, which is usually EEVDF \parencite[EEVDF Scheduler]{KernelSchedManual} in the Linux environment. Processes are usually scheduled to run on kernel threads, and the kernel threads are scheduled to have equal time allotted on the CPU. Multiple threads within a single process share the same memory space, making inter-thread communication faster but also introducing synchronisation challenges.

Processes maintain several important attributes beyond their PID, including a Parent Process ID (PPID) that identifies which process created them, a user ID (UID) and group ID (GID) that determine privileges, and various resource limits. When a parent process terminates before its children, those child processes are typically reparented to the init process (PID 1), becoming orphaned processes. Conversely, when a child terminates but the parent has not yet read its exit status, it becomes a zombie process, retaining an entry in the process table until the parent collects its status \parencite{GnuBashManual}.

\subsubsection{Pipelines and Process Communication}
\label{PipelinesAndProcessComm}

A pipeline in the context of the Linux commandline refers to the mechanism of connecting the output of one process to the input of another process using the \texttt{|} operator \parencite{GFGPiping}. This is not related to the concept of pipelines in the context of data processing or machine learning, but it is a powerful tool for chaining commands together and creating complex workflows. An example of this would be the command \texttt{ps aux | grep python}, which would list all the processes running on the system and then filter the results to show only those that contain the word "python". This allows for a user to easily search through the processes and find specific ones that they are interested in.

Pipelines are implemented using anonymous pipes, which are unidirectional data channels managed by the kernel. When a shell creates a pipeline, it forks multiple child processes and connects their standard output and input streams. The processes in a pipeline execute concurrently, with data flowing between them as it becomes available. All processes in a pipeline share the same process group ID (PGID), which allows them to be managed collectively.

Beyond anonymous pipes, Linux provides several other inter-process communication (IPC) mechanisms. Named pipes (FIFOs) persist in the filesystem and can be used by unrelated processes. Shared memory allows multiple processes to access the same memory regions for high-performance data exchange. Message queues provide structured communication channels. Sockets enable communication both locally and across networks. Understanding these mechanisms is crucial because they represent potential channels through which AI agents might exchange information or detect each other's presence \parencite[3.2.2]{GnuBashManual}.

\newpage
\subsubsection{Job Control and Management}
\label{JobControl}

By cross referencing \textcite{GnuBashManual,GFGJobSPEC,GFGJobScheduling,GFGCronCommand}, jobs are a concept in the Linux commandline manuals that refer to processes and pipelines that are either started or scheduled to run by the user in a shell session or from the terminal. They are given a process ID (PID) on creation, regardless of whether they are multiple commands tied together in a pipeline or a single command. They can be run in the foreground or background, and they can be managed using built-in commands such as \texttt{jobs}, \texttt{fg}, \texttt{bg}, and \texttt{kill}.

Foreground jobs receive keyboard input and prevent the shell from accepting new commands until they complete or are suspended. Background jobs, launched by appending \texttt{\&} to a command, run concurrently with the shell, allowing the user to continue issuing commands. The \texttt{Ctrl-Z} key combination suspends a foreground job, which can then be resumed in the background with \texttt{bg} or returned to the foreground with \texttt{fg}.

Job control is mediated through signals, which are asynchronous notifications sent to processes. Common signals include SIGTERM (polite termination request), SIGKILL (immediate forced termination), SIGSTOP (suspend execution), and SIGCONT (resume execution). The \texttt{kill} command, despite its name, is the primary interface for sending arbitrary signals to processes \parencite[7.1-7.3]{GnuBashManual}.

Using cron, a user can also schedule jobs to run at specific times or intervals, which can be useful for automating tasks and running processes without user intervention \parencite{GFGCronCommand}. The \texttt{crontab} command manages per-user cron schedules, while system-wide cron jobs are defined in \texttt{/etc/crontab} and \texttt{/etc/cron.d/}. Each entry in a crontab file follows a five-field time specification of the form \texttt{* * * * * command}, where the fields represent minute (0--59), hour (0--23), day of month (1--31), month (1--12), and day of week (0--6), respectively \parencite{GFGCronCommand}. The \texttt{at} and \texttt{batch} commands provide one-time scheduled execution. These scheduling mechanisms could allow an AI agent to persist its presence or execute actions at predetermined times.

Where processes usually differ from jobs is that they can be triggered and scheduled by the operating system without user involvement. Jobs are generally more manageable for users, as they are directly tied to the shell and can be easily manipulated using built-in commands. Processes, on the other hand, can be more complex and may require more advanced tools for management and monitoring. Due to their nature as well, processes can be more anonymous and harder to track as well, seeing as they can be triggered by other processes and may not have a direct association with a user or shell session.

\subsubsection{Process Monitoring and Discovery}
\label{ProcessMonitoringDiscovery}

Commands such as \texttt{ps}, \texttt{top}, \texttt{htop} and \texttt{pgrep} can be used to monitor and search through processes, whilst commands such as \texttt{pkill}, \texttt{killall} in addition to \texttt{kill} all allow for termination of processes \parencite{LinuxManPages}.

The \texttt{ps} command provides a snapshot of current processes. With options like \texttt{aux}, it displays all processes system-wide with detailed information including CPU and memory usage, execution time, and the command line used to launch each process. The \texttt{-ef} option provides a full-format listing showing parent-child relationships through PPID values. For AI agents, examining process command lines and parent relationships could reveal the presence of competing agents or suspicious activity.

The \texttt{top} and \texttt{htop} commands provide real-time, continuously updating views of system processes, sorted by various metrics such as CPU or memory consumption. These interactive tools highlight processes that are consuming unusual amounts of resources, which could indicate malicious behaviour or intensive computational tasks associated with an AI agent.

The \texttt{pgrep} command searches for processes by name or other attributes, returning only the PIDs. Its counterpart \texttt{pkill} sends signals to matching processes. The \texttt{pidof} command finds PIDs by exact command name. The \texttt{pstree} command displays the process hierarchy as a tree, making parent-child relationships visually apparent.

More sophisticated process introspection is available through the \texttt{/proc} pseudo-filesystem, where each process has a directory named by its PID containing files with detailed information: \texttt{/proc/[pid]/cmdline} shows the command line, \texttt{/proc/[pid]/exe} is a symlink to the executable, \texttt{/proc/[pid]/environ} contains environment variables, and \texttt{/proc/[pid]/fd/} lists open file descriptors. An AI agent with filesystem access could extensively analyse other processes through this interface \parencite{LinuxManPages}.

\subsubsection{Process Control and Termination}
\label{ProcessControlTermination}

Beyond the basic \texttt{kill} command, Linux provides several mechanisms for process termination. The \texttt{killall} command sends signals to all processes matching a given name, while \texttt{pkill} offers more flexible pattern matching and attribute-based selection. The \texttt{timeout} command can be used to automatically terminate a process if it exceeds a specified execution time.

Process resource limits can be queried and modified using \texttt{ulimit} (for the current shell and its children) or \texttt{prlimit} (for arbitrary processes). These limits cover maximum file sizes, number of open files, memory usage, CPU time, and other resources. An AI agent could potentially use these mechanisms to constrain a competing agent's capabilities or to detect constraints imposed upon itself.

The \texttt{strace} command traces system calls made by a process, revealing every interaction with the kernel including file operations, network activity, and process management. Similarly, \texttt{ltrace} traces library calls. These powerful debugging tools could allow one AI agent to monitor another's behaviour in detail, though they require appropriate privileges \parencite{LinuxManPages}.

\subsubsection{Process Scheduling and Priority Management}
\label{ProcessScheduling}

In addition to this, the EEVDF scheduler \parencite[EEVDF]{KernelSchedManual} allows for scheduling affinity using the \texttt{sched\_setattr()} system call, which the user can call on directly on starting a process using the \texttt{nice} command, or after the process has started with \texttt{renice}. This allows for a user to give a process higher or lower priority on the CPU, which can make some processes run with more CPU time than others and vice versa. This can be used to give a process more resources to run, or to starve a process of resources and make it run slower.

Process priorities in Linux are expressed as "nice" values ranging from -20 (highest priority) to 19 (lowest priority), with 0 being the default. Despite the name, higher nice values are "nicer" to other processes by yielding more CPU time to them. Only privileged users can decrease nice values (increase priority), though any user can increase nice values (decrease priority) for their own processes.

Beyond nice values, Linux supports multiple scheduling classes. The default SCHED\_OTHER (or SCHED\_NORMAL) uses the EEVDF algorithm for time-sharing. Real-time scheduling classes (SCHED\_FIFO and SCHED\_RR) provide deterministic scheduling for time-critical applications but require elevated privileges. The \texttt{chrt} command can query and modify scheduling policies and priorities.

CPU affinity, controlled by \texttt{taskset}, restricts a process to execute only on specified CPU cores. This can be used for performance optimization but could also be weaponized to isolate or constrain competing processes. The \texttt{cgroups} (control groups) mechanism provides comprehensive resource management, allowing fine-grained control over CPU time, memory allocation, I/O bandwidth, and other resources for groups of processes \parencite{KernelSchedManual}.

\subsubsection{Process Dependencies and Synchronisation}
\label{ProcessDependencies}

In addition to these, there are also commands such as \texttt{wait} which can be used to pause the execution of a process until another process has finished, and \texttt{fuser} which can be used to identify which processes are using a specific file or socket. This can be useful for identifying which processes are interacting with each other and potentially interfering with each other.

The \texttt{wait} command is primarily used in shell scripts to synchronise parent processes with their children. Without \texttt{wait}, a parent might terminate before its children complete, or might not know whether child processes succeeded or failed. The command can wait for specific PIDs or for all background jobs \parencite[3.2]{GnuBashManual}.

The \texttt{fuser} command identifies processes using specific files or sockets, returning PIDs and access types (read, write, execute, or root directory). This is particularly valuable for detecting filesystem-based interactions between processes. The related \texttt{lsof} (list open files) command provides comprehensive information about all files opened by processes, including regular files, directories, network sockets, pipes, and devices. Since everything in Unix-like systems is treated as a file, \texttt{lsof} offers a powerful view into process behaviour and inter-process relationships \parencite{LinuxManPages}.

File locking mechanisms, accessible through the \texttt{flock} command or programmatically through system calls, allow processes to coordinate access to shared resources. An AI agent could use file locks both for legitimate synchronisation and for detecting or blocking competing agents attempting to access the same resources \parencite{KernelFSLocking}.

\subsection{User Management and Groups}
\label{UserManagement}

Linux is a multi-user operating system, meaning that several accounts can be active on the same system simultaneously, each with its own set of privileges and resource limits. Every process runs under the identity of a specific user, identified by a numeric user ID (UID) and one or more group IDs (GIDs). These identities determine what files the process can read, write, or execute, and which system operations it is permitted to perform \parencite{KernelCredentials}. Beyond the real UID and GID, each process also carries an effective UID (EUID) and effective GID (EGID), which the kernel uses as the subjective context when evaluating whether an operation is permitted; these may differ from the real identifiers when a process has changed its privileges through system calls or privilege escalation mechanisms \parencite{KernelCredentials}.

The \texttt{id} command reports the current user's UID, primary GID, and all supplementary group memberships. The \texttt{whoami} command returns just the username, while \texttt{who} and \texttt{w} list all users currently logged into the system along with their terminal sessions and recent activity. These commands are relevant to the AESM because an agent could use them to enumerate other active sessions and infer whether a rival agent is operating under a different user account \parencite{LinuxManPages}.

User and group accounts are managed through \texttt{useradd}, \texttt{usermod}, and \texttt{userdel} for users, and \texttt{groupadd}, \texttt{groupmod}, and \texttt{groupdel} for groups. Passwords are set or changed with \texttt{passwd}. These administrative commands typically require root privileges and are relevant to understanding the boundaries within which agents are constrained \parencite{KernelCredentials}.

Privilege escalation is achieved through \texttt{su} (substitute user), which opens a new shell under a different user's identity after supplying that user's password, and \texttt{sudo} (superuser do), which allows a permitted user to execute a single command as another user, most commonly root, as configured in \texttt{/etc/sudoers} \parencite{LinuxManPages}. The distinction between unprivileged and root-level operation has significant implications for the AESM, as discussed further in \autoref{PrivilegeLevel}: an agent running as root faces almost none of the access restrictions that constrain an unprivileged agent, and can freely terminate processes or modify files owned by any other user.